% Generated draft: requires mathematical and source review.
\documentclass[11pt]{article}
\usepackage[margin=25mm]{geometry}
\usepackage{fontspec}
\setmainfont{Latin Modern Roman}
\usepackage{amsmath,amssymb,mathtools,unicode-math}
\setmathfont{Latin Modern Math}
\usepackage{enumitem,needspace,xurl,hyperref}
\setlength{\parindent}{0pt}
\setlength{\parskip}{5pt}
\setlength{\emergencystretch}{3em}
\newcommand{\N}{\mathbb{N}}
\newcommand{\Z}{\mathbb{Z}}
\newcommand{\Q}{\mathbb{Q}}
\newcommand{\R}{\mathbb{R}}
\newcommand{\C}{\mathbb{C}}
\newcommand{\F}{\mathbb{F}}
\newcommand{\A}{\mathbb{A}}
\newcommand{\PP}{\mathbb{P}}
\newcommand{\E}{\mathbb{E}}
\newcommand{\eps}{\varepsilon}
\newcommand{\dd}{\,\mathrm{d}}
\begin{document}
\section*{UnsolvedMath ID: 20003175}
\textbf{Source problem code:} catalogue entry
\subsection{Definitions and mathematical statement}
The source asks whether Ricci flow can upgrade the Hsiang--Kleiner homeomorphism classification of positively curved four-manifolds with symmetry to a diffeomorphism classification. The trailing heading ``8. Miscellaneous Problems'' is not part of the question. In the original oriented setting, let \((M^4,g_0)\) be closed, connected, oriented, with strictly positive sectional curvature and an effective isometric \(S^1\)-action. The intended standard smooth types are \(S^4\) and \(\C\PP^2\), with the orientation of \(\C\PP^2\) reversed when required.

The supplied aid says the diffeomorphism conclusion is already known by Grove--Wilking, Geometry \& Topology 18 (2014), Theorem A: a closed positively curved four-manifold with effective isometric circle symmetry is equivariantly diffeomorphic to a linear action on \(S^4\), \(\R\PP^4\), or \(\C\PP^2\). In the oriented Hsiang--Kleiner setting the \(\R\PP^4\) case is excluded. This supplied theorem resolves the smooth classification, but its proof is not a four-dimensional curvature-preserving Ricci-flow deformation. The method-specific request remains distinct from the classification result.

The aid reports a persistence lemma. For a smooth Ricci flow on the closed fixed manifold, either
\[
\partial_tg=-2\operatorname{Ric}(g)
\quad\text{or}\quad
\partial_tg=-2\operatorname{Ric}(g)+c(t)g
\]
with a spatially constant normalization, the original \(S^1\)-action remains effective and isometric throughout its smooth existence interval. Its fixed sets, isotropy groups, slice representations/weights and legally weighted orbit data remain those of the same underlying smooth action. Thus such a flow does not simplify a weighted exceptional circle merely by changing the metric. This assertion does not include passage through singularities, surgery or a change of underlying action.

In the original oriented hypotheses, the aid also reports that a smooth extension of the circle action to an effective \(T^2\)-action, together with the Hsiang--Kleiner bound \(\chi(M)\le3\), forces the standard smooth structure \(S^4\) or oriented \(\pm\C\PP^2\). Producing that extension is an additional step, not a consequence of symmetry persistence alone. A flow-based strategy would need suitable canonical convergence/surgery or control of the knotting of weighted exceptional orbit data in the three-dimensional quotient. The source does not provide such a strategy as a theorem.

Finally the aid cites Bettiol--Krishnan examples showing that positive sectional curvature need not be preserved by four-dimensional Ricci flow. Their full symmetry and initial-metric hypotheses are not supplied here, so no stronger counterstatement for every Hsiang--Kleiner flow is invented. The examples rule out assuming general sectional-curvature preservation as an automatic step.
\subsection{Short English statement}
The supplied literature already gives the smooth classification, while the requested Ricci-flow proof is a separate method question. Smooth flow preserves the original circle action and its orbit data, and positive sectional curvature cannot be assumed to persist without further hypotheses.
\subsection{Sources}
\begin{itemize}
\item Problem 50, AIM workshop \emph{Manifolds with nonnegative sectional curvature}, \url{https://aimath.org/WWN/nnsectcurvature/nnsectcurvature.pdf}.
\item Grove--Wilking, Geometry \& Topology 18 (2014), Theorem A, as supplied for the equivariant diffeomorphism classification. The record does not give its full title or page range.
\item The aid names the Hsiang--Kleiner theorem and Euler bound, and Bettiol--Krishnan's positive-curvature Ricci-flow counterexamples, without complete bibliographic entries or full hypotheses.
\item Local UnsolvedMath registry, permanent ID 20003175, exact code \texttt{catalogue entry}; canonical location \texttt{aim-topology-notes.json}, notes[262], dated 14 August 2026; statement labeled exact, registry status solved and research classification solved in literature.
\item Research attempt 2: partial result, problem solved at the run, rigor label proved, novelty confidence low. The persistence and torus-extension reductions are machine-generated research aids, not peer reviewed; independent verification is required before citation.
\end{itemize}
\end{document}
